\section{Discussion}\label{sec:discussion}

The theorems in Section~\ref{sec:theorems} preserve the oracle when the local laws hold, and the experiments in Sections~\ref{sec:markov-walkthrough}, \ref{sec:hll-parity}, and~\ref{sec:manifesto-llm} exercise that preservation on three progressively richer oracles. The resulting trade-off, in plain language: an audited structural term, a fixed calibration residual, and a sampling envelope that shrinks with the total label count.

\subsection{The Finite-Sample Gap Decomposition}\label{sec:discussion-tradeoff}

Under the local laws, the finite-sample gap between a preference
objective on tree summaries and the corresponding full-document
objective decomposes into three non-negative pieces:
\begin{equation}\label{eq:discussion-bound}
|G|
\;\le\;
\underbrace{\hat{\Delta}_R}_{\text{empirical IPW estimate}}
\;+\;
\underbrace{B_{\mathrm{cal}}}_{\text{fixed residual}}
\;+\;
\underbrace{\text{SE}(n_{\mathrm{global}} + n_{\mathrm{local}})}_{\text{shrinking envelope}}.
\end{equation}
This is the decomposition the main gap-bound theorem gives on any audit-produced certificate (Appendix~\ref{app:proof-artifacts}). The three pieces behave differently, and each one reads most cleanly under one side of the projection/structured-error iff.

The empirical term $\hat{\Delta}_R$ (read under the \emph{structural-error} framing) is an IPW union bound over sampled leaves, merges, and idempotence rounds; equivalently, it is the audit's empirical estimate of how far the approximation error fails to vanish on the C1/C2/C3 subspace. It is zero when the audit reports zero violations and grows linearly with the realized violation rate. Section~\ref{sec:fno-primer} and Appendix~\ref{app:approx-via-neural-operators} give the approximation-theoretic counterpart: a practitioner chooses an upstream operator tolerance $\varepsilon$ on compact realized calls and then converts that tolerance, through explicit transfer moduli, into audited leaf/merge/idempotence budgets. Accordingly, $\hat{\Delta}_R$ is the audited manifestation of a representation-design target, not a term set by fiat.

The fixed residual $B_{\mathrm{cal}}$ (read under the \emph{projection} framing) is the irreducible projection gap between the judge's C1/C2/C3 subspace and the oracle's: the gap between the judge used during training and the task oracle the researcher ultimately cares about. It is a function of the calibration set and the confidence level only, so it does not depend on how many downstream labels the practitioner subsequently collects: if the judge is the oracle, $B_{\mathrm{cal}} = 0$; otherwise it is the irreducible residual of the measurement instrument.

The envelope $\text{SE}(n_{\mathrm{global}} + n_{\mathrm{local}})$ is the part that pays for labels; it is framing-agnostic, because it is sampling noise on whichever of the two framings the audit is measuring. The idealized envelope $C / \sqrt{n+1}$ is non-negative, bounded above by the structural constant $C$, and strictly decreasing in $n$ whenever $C > 0$; the separate-monotonicity statement extends this to two label types, so holding one kind of label fixed, adding the other still tightens the envelope. Concretely, $n_{\mathrm{global}}$ counts full-document labels and $n_{\mathrm{local}}$ counts sampled audit nodes; both feed the same envelope.

The function-space view also clarifies what the training regularizer is
doing. The per-law-weighted generalization of
Equation~\eqref{eq:oracle-projection-objective}
(Appendix~\ref{app:oracle-plus-projection}) has one term fitting the
oracle and a $\lambda$-weighted local-law term pushing the realized
operator toward the local-law-constrained subspace. The endpoint theorems make the
interpretation literal: $\lambda=0$ is pure oracle risk, and $\lambda=1$
is pure projection penalty. Appendix~\ref{app:projection-interpretation}
spells out the corresponding iff: treating the local-law weights as a
projection and treating the approximation error as structured by C1/C2/C3
are the same formal assumption. That is why the same language applies to
the examples above: HLL starts inside the subspace, the Markov tree learns
an ordered boundary-sensitive member of it, and the manifesto pipeline
uses audits to estimate how close the LLM operator lands to it.

Long context windows do not remove that issue; they change where the
engineering pressure appears. A million-token window may make truncation
less common, but it does not certify that a single read made the same
rubric decision an expert would have made after tracking evidence across
the whole document. The local architecture gives a different claim: if
small leaves and pairwise merges preserve the oracle-relevant evidence,
then the root measurement is stable under divide-and-conquer and the
audit has localized places to look when it is not. The manifesto results
are meant in exactly that sense, not as a claim that context length is
irrelevant.

\subsection{Global and Local Labels}

The global/local separation is exactly what the budget-split
experiment in Section~\ref{sec:markov-walkthrough} demonstrates.
Replacing some fraction of root labels by a proportionate amount of
local labels (at matched token mass) leaves root accuracy essentially
unchanged all the way down to about $40\%$ root share, because the
envelope responds to the \emph{total} label count, not to the ratio
of its two components. The global/local decomposition theorem referenced
above is the formal statement that backs this empirical observation.
Leaf-size invariance is what makes the same trade meaningful in
long-document settings: when changing the leaf grid leaves the root
estimate stable, each smaller leaf becomes another less expensive site
for sampling and supervision rather than a change in the estimand.

\subsection{Scope}

Equation~\eqref{eq:discussion-bound} is a conditional guarantee: it
assumes the local laws hold on the realized tree, which
Section~\ref{sec:estimation}'s audit checks. Tasks whose oracle does
not decompose into a locally checkable sufficient statistic (e.g.\
rhetorical-tone or discourse-level semantics in which replacing a
span with a locally equivalent one changes how neighboring text is
interpreted) will show high C1/C3 violation rates under the audit,
and the bound will honestly report that the tree's root prediction is
no better than a flat baseline on such tasks. The framework does not
``always'' apply; it applies \emph{when the audit says it applies}.

\subsection{Applications in Social-Scientific Measurement}

The experiments span three points on the classical-to-learned
gradient. At the classical limit (Section~\ref{sec:hll-parity}), the
oracle is a distinct-count query, the merge is algebraically exact,
$B_{\mathrm{cal}} = 0$, and the envelope collapses: the tree is
byte-identical to the flat reference. On the Markov changepoint task
(Section~\ref{sec:markov-walkthrough}), the oracle is additive with a
boundary correction; the tree pays a small estimation envelope but
matches or improves on the flat baseline across the supervision
budget. On Manifesto-Project RILE scoring (Section~\ref{sec:manifesto-llm})
the oracle is rubric-driven, $B_{\mathrm{cal}}$ is the judge/expert
gap, and the envelope shrinks with either pile of labels a
practitioner can supply.

Many social-scientific measurement problems share the Manifesto
shape: rubrics with locally codable categories, span-grounded
expert judgments, and aggregates that decompose across sections,
chapters, or acts. Comparative-politics party-position scoring
\citep{BenoitEtAl2025,ManifestoProject2025a}, policy-agenda coding,
clinical narrative review, and legal-corpus analysis all fit this
description. The practical value of the framework in these settings
is that Equation~\eqref{eq:discussion-bound} turns long-document
annotation into an audit-budget problem: a fixed calibration set
sets $B_{\mathrm{cal}}$ once, and any combination of global and
local labels the researcher can collect tightens the envelope along
the monotone curve $C/\sqrt{n+1}$.
